% ---------------------------------------------------------------------------------------------
% Section 4: Expected Findings and Evaluation Plan (#14). Paste after Section 3.
% Our project is an empirical study (RQ1, RQ2) plus a tool (RQ3), so this section covers both:
% hypotheses and implications for the study, and benchmark / baselines / metrics for ReliGuard.
% ---------------------------------------------------------------------------------------------

\section{Expected Findings and Evaluation Plan}
\label{sec:expected}

\paragraph{Benchmark and metrics.}
The benchmark is Version~0, the eight change requests, and the 12 property checks; the 72 versions
the agents produce are our data. For each version we record the functional pass rate (the agent's
own tests and our nine functional tests), which properties regress, and every value they measured.
For RQ1 we report the fraction of versions with at least one regression, the number of properties
broken per version, and how both develop from change~1 to change~8, alongside the functional pass
rate. For RQ2 we cross-tabulate regressions by change and by fault type, and two authors
independently label each regression's root cause from the diff (for example removed timeout,
unbounded retry, deadline bypassed, idempotency key lost, compensation skipped), resolving
disagreements by discussion. For RQ3 the unit is one regression. The same agent repairs it twice
from the same starting version, once with \textsc{ReliGuard}'s evidence and once with a baseline that gives
only the failing check's name and raw output, the feedback used by prior repair
work~\cite{chen2024selfdebug,xia2024chatrepair}. We measure \emph{repair success} (the original
property passes again), \emph{clean repair} (the whole suite passes, D4), \emph{collateral
regressions} (properties the repair newly breaks), and \emph{functional preservation} (the requested
feature still works), plus turns and cost.

\paragraph{H1: functional tests miss most reliability regressions.}
We expect a large gap between the two pass rates: agents' own tests will pass in nearly every
version, while a substantial share of versions, in the tens of percent, will break at least one
property. The pilot already shows this pattern (47/47 tests passing, 3 of 10 properties broken).
We also expect regressions to accumulate. Once a change removes a guarantee, later changes have
no test that would restore it, so the number of broken properties should not decrease along the
sequence. If confirmed, this means ``all tests pass'' is a weak acceptance criterion for
agent-maintained services, and that benchmarks that evaluate agents over long change sequences
with functional tests alone~\cite{deng2026evoclaw,le2025sweevo} overestimate how safe those
agents' changes are.

\paragraph{H2: the change determines the property.}
We expect regressions to concentrate in the changes that touch the call path (payment retries, the
payment timeout, the notification service on the checkout path, the database table) and in the
properties about time and repetition (P1--P4), more than in idempotency (P5, P6), which the
original code makes explicit. We also expect a common cause: agents add a new limit (a retry window,
a timeout) without connecting it to the existing end-to-end deadline, as change~02 did. If this
holds, reviewers and tools can focus fault injection on a small set of change types instead of
running every check everywhere, and agent instructions can name the specific invariants that
tend to break.

\paragraph{H3: structured evidence produces cleaner repairs.}
We expect both feedback conditions to restore the failing property often, because the failing test
name already points to the problem. The difference should appear in clean repair and collateral
regressions. With only a raw log, an agent tends to patch the symptom (for example, lowering one
timeout so a test passes) and can break another property. With the injected fault, the measured
values, and the diff that caused the regression, it can restore the original mechanism. If
confirmed, this supports building fault-injection feedback into agent workflows as a standard step
after each change, and it shows that the kind of feedback matters, not just the amount.

\paragraph{Threats.}
Our results come from one application that we wrote ourselves, and from three agents at fixed
versions. We mitigate this by modeling Version~0 on a widely used reference
application~\cite{onlineboutique}, releasing every prompt, check, and transcript, and validating
the checks for stability and sensitivity (Section~\ref{sec:props}). The findings describe how these
agents behave on this benchmark; showing that they generalize to other systems is future work.
